% --- Chapter appendices, numbered <chapter>.A, .B, ... -----------------------
\renewcommand{\thesection}{\thechapter.\Alph{section}}
\setcounter{section}{0}

% Appendix A --- the birth--plague model, reduced to a bare specification.
% Was \S1.5 of the body.  Cut: the motivation, the rate diagram, the
% discussion of what the coupling does, and the open questions.  One figure
% retained.
% The NEEDS-AUTHOR question on the vitality variant stands: the source
% (09/09:11) specifies dV/dt = \xi - p X_t and then uses V in none of the
% transition rates.  Written here as corruption at \delta X_t / V(t), by
% consistency with the catastrophe channel of \cref{var:eq:rrrr}, which is the
% model the variant says it is borrowing from.  The alternative reading is
% \delta X_t^2 / V(t), under which the loss channel is quadratic in X and the
% collapse is sharper.  Author to confirm; nothing else depends on it.

\section{The birth--plague model}
\label{var:app:plague}

Where \cref{var:sec:rrrr} drives the loss rate by a potential carried
alongside a single count, this model drives it by a second count, so that
what accelerates the loss of one component is the size of the other. It is
recorded here as a specification; it has been simulated but not analysed.

Let $X_t$ count uninfected individuals and $Y_t$ the infected. In a vanishing
interval,
\begin{align}
\pr{\Delta X_t=1,\ \Delta Y_t=0\mid X_t,Y_t}
  &=\lambda X_t\,\Delta t+o(\Delta t),
  \label{var:eq:plague}\\
\pr{\Delta X_t=-1,\ \Delta Y_t=1\mid X_t,Y_t}
  &=\bigl(\delta+\chi Y_t\bigr)X_t\,\Delta t+o(\Delta t),\notag\\
\pr{\Delta X_t=0,\ \Delta Y_t=-1\mid X_t,Y_t}
  &=\mu Y_t\,\Delta t+o(\Delta t),\notag\\
\pr{\Delta X_t=\Delta Y_t=0\mid X_t,Y_t}
  &=1-\bigl[(\lambda+\delta+\chi Y_t)X_t+\mu Y_t\bigr]\Delta t+o(\Delta t),
  \notag
\end{align}
with $\lambda$ the birth rate, $\delta$ the rate of spontaneous corruption,
$\chi$ the contact-infection coefficient and $\mu$ the clearance rate of
infected individuals.

The middle channel is written once and carries two mechanisms superposed. An
uninfected individual is lost to corruption spontaneously at the constant
per-capita rate $\delta$, and to infection by contact at rate $\chi$ per
uninfected--infected pair. Both produce the same jump, so the only thing that
distinguishes them is how their rates scale: the first is linear in $X_t$,
the second bilinear in $X_tY_t$. Only that channel carries $Y_t$, but the
coupling is two-way, since every infection removes one $X$ and adds one $Y$.

\paragraph{The vitality-coupled variant.} Carry a vitality $V(t)$ alongside
the pair, replenished at a constant rate and drawn down in proportion to the
population,
\begin{equation}
\ddt{V}=\phi-c\,X_t ,
\label{var:eq:vitalityV}
\end{equation}
and divide the spontaneous part of the loss channel by it:
\begin{equation}
\pr{\Delta X_t=-1,\ \Delta Y_t=1\mid X_t,Y_t}
=\lrb{\frac{\delta}{V(t)}+\chi Y_t}X_t\,\Delta t+o(\Delta t),
\label{var:eq:vitality}
\end{equation}
the other three channels being unchanged. The reason for wanting the variant
is that \cref{var:eq:plague} has no mechanism for a collapse: its infected
count rises when $X_t$ is large and falls when it is small, which damps
rather than amplifies. Dividing corruption by a vitality the population
itself depletes supplies the missing amplification, since a population which
has grown large drives $V$ down and so raises its own loss rate --- the
ingredient \cref{var:sec:rrrr} used to close its cycle. Whether the
combination gives a sustained cycle or a single collapse followed by a
low-level steady state is the question the variant exists to ask. It has not
been simulated.

\begin{figure}[!htbp]
\centering
\includegraphics[width=0.9\textwidth]{fig09_plague_paths/fig09}
\caption{Two realisations of \cref{var:eq:plague} at
$(\lambda,\delta,\chi,\mu)=(1,\,0.1,\,0.02,\,10)$: uninfected count $X_t$ in
blue, infected count $Y_t$ in vermillion. The clearance rate is large
relative to the contact rate, so infected individuals never accumulate and
the epidemic acts as a brake on $X_t$ rather than as a collapse of it. Set
beside \cref{var:fig:vbc-many}, the fluctuations never reset a phase.}
\label{var:fig:bpplots}
\end{figure}

\FloatBarrier
